% Chapter: The Refinement-Fibered C⁴ Compiler
% What Cohomology Adds Beyond Cubical Type Theory — In Practice
%
% This chapter makes the concrete case that the cohomological
% component of C⁴ is not mere theoretical decoration but provides
% genuine gains in TRACTABILITY and EXPRESSIVITY for Python verification.
% It introduces the refinement site (replacing duck-type fibers),
% proves tractability theorems, and gives the full VC compilation
% algebra for the c4_compiler.

\chapter{The Refinement-Fibered $\mathrm{C}^4$ Compiler}
\label{ch:c4-compiler}

\begin{quote}
\emph{``The purpose of abstraction is not to be vague, but to create
a new semantic level in which one can be absolutely precise.''}
--- Edsger Dijkstra
\end{quote}


% ═══════════════════════════════════════════════════════════════
\section{Introduction: The Verification Bottleneck}
\label{sec:compiler-intro}

The preceding chapters developed $\mathrm{C}^4$ as a calculus:
site-decorated universes, cubical paths, the descent rule,
Mathlib import.  This chapter addresses the \emph{engineering}
question: how does this calculus compile to something a machine
can check?  Specifically, how does one take a proof term built
from the constructors of Chapter~\ref{ch:c4-calculus} and produce
a finite set of \textbf{verification conditions} (VCs) that
Z3 can decide?

The bottleneck in automated program verification is not expressivity
--- dependent type theories are arbitrarily expressive.  The
bottleneck is \emph{tractability}: most interesting properties of
most interesting programs lie outside the decidable fragment of any
fixed logic.  A verification system is useful not when it can
\emph{express} all properties, but when it can \emph{decide} a
large fraction of properties that arise in practice.

The thesis of this chapter is:

\begin{quote}
\textbf{Thesis.}
The cohomological component of $\mathrm{C}^4$ provides a systematic
mechanism for converting \emph{intractable global} verification
problems into \emph{tractable local} ones, with decidable gluing
conditions.  This is a genuine gain in tractability that cubical
type theory alone does not provide.
\end{quote}

We develop this thesis through three contributions:

\begin{enumerate}
\item \textbf{The refinement site}
  (Section~\ref{sec:refinement-site}): a new site category for
  $\mathrm{C}^4$ where the objects are \emph{refinement predicates}
  $\{x : A \mid \varphi(x)\}$ rather than the coarse duck-type
  fibers of Chapter~\ref{ch:c4-descent}.  This site is richer,
  finer, and Z3-decidable.

\item \textbf{Tractability theorems}
  (Section~\ref{sec:what-cohomology-adds}): formal results showing
  that cohomological descent reduces proof complexity from
  exponential to linear in the number of branches, with each
  local obligation in a decidable fragment.

\item \textbf{The VC compiler}
  (Sections~\ref{sec:vc-compiler}--\ref{sec:worked-examples}):
  the complete compilation algebra mapping each proof rule to
  explicit Z3 obligations, with cubical face-map checks, trust
  provenance, and F*-style annotation binding.
\end{enumerate}


% ═══════════════════════════════════════════════════════════════
\section{The Refinement Site}
\label{sec:refinement-site}

Chapter~\ref{ch:c4-descent} defined the CCTT site $\site_\mathrm{CCTT}$
with objects being coarse duck-type fibers: $\mathsf{TInt}$,
$\mathsf{TFloat}$, $\mathsf{TStr}$, etc.  This site was convenient
for the equivalence-checking problem but is too coarse for
\emph{verification}: two functions can agree on all duck-type fibers
yet violate a refinement property (e.g., non-negativity, sortedness,
non-nullity).

We now define a finer site whose objects capture the actual
control-flow structure of Python programs.

\begin{definition}[Refinement Predicate]
\label{def:refinement-pred}
A \textbf{refinement predicate} $\varphi$ is a formula in the
quantifier-free fragment of the theory of linear integer arithmetic
extended with uninterpreted functions (QF\_UFLIA), optionally
involving:
\begin{enumerate}
\item Arithmetic comparisons: $x > 0$, $\mathit{len}(xs) \geq n$.
\item Type tests: $\mathit{isinstance}(x, \tau)$.
\item Nullity tests: $x\ \mathit{is}\ \mathsf{None}$,
  $x\ \mathit{is\ not}\ \mathsf{None}$.
\item Boolean connectives: $\varphi_1 \wedge \varphi_2$,
  $\varphi_1 \vee \varphi_2$, $\neg \varphi$.
\item Library predicates (as uninterpreted):
  $\mathit{sympy.is\_polynomial}(e)$, $\mathit{sorted}(xs)$.
\end{enumerate}
A refinement predicate is \textbf{Z3-decidable} if it lies in a
decidable theory for which Z3 has a decision procedure (e.g.,
QF\_LIA, QF\_UFLIA, QF\_BV).  A refinement predicate involving
library predicates (item~5) is \textbf{library-decidable}: decidable
given library axioms, but not in pure Z3.
\end{definition}

\begin{definition}[The Refinement Site $\site_\mathrm{Ref}$]
\label{def:refinement-site}
The \textbf{refinement site} $\site_\mathrm{Ref}$ is a category
with a Grothendieck topology:
\begin{enumerate}
\item \textbf{Objects:} refinement predicates $\varphi(x)$ over a
  fixed variable context $\Gamma = x_1 : \tau_1, \ldots, x_k : \tau_k$.
  Each object represents the set of inputs satisfying $\varphi$.

\item \textbf{Morphisms:} there is a (unique) morphism
  $\varphi \to \psi$ whenever $\varphi(x) \Rightarrow \psi(x)$ is
  valid.  Equivalently, $\{x \mid \varphi\} \subseteq \{x \mid \psi\}$.
  This is decidable by Z3 for Z3-decidable predicates.

\item \textbf{Covering families:} a family $\{\varphi_i\}_{i \in I}$
  covers $\psi$ if $\psi \Rightarrow \bigvee_i \varphi_i$ is valid.
  The \textbf{standard cover} is the family of branch predicates
  extracted from a function's control flow.

\item \textbf{Terminal object:} $\top$ (the predicate ``true''),
  representing the entire input domain.
\end{enumerate}
\end{definition}

\begin{remark}[Comparison with the Duck-Type Site]
\label{rem:duck-vs-refinement}
The refinement site is a \emph{strict refinement} of the duck-type
site: every duck-type fiber $\mathsf{TInt}$ embeds as the refinement
predicate $\mathit{isinstance}(x, \mathsf{int})$.  But the refinement
site is vastly richer:
\begin{itemize}
\item The duck-type site has $\sim 10$ objects.
  The refinement site has infinitely many, though for any fixed
  function only finitely many are relevant (those appearing in the
  control flow).
\item The duck-type site captures \emph{type structure}.
  The refinement site captures \emph{value structure}: which
  inputs are positive, which lists are sorted, which dicts
  have a certain key.
\item The duck-type site has morphisms only from Python's subtype
  hierarchy ($\mathsf{bool} \hookrightarrow \mathsf{int}
  \hookrightarrow \mathsf{float}$).  The refinement site has
  morphisms from \emph{logical implication}: $x > 5 \Rightarrow x > 0$
  is a morphism, letting proofs about positive numbers specialize
  to proofs about numbers greater than~5.
\end{itemize}
\end{remark}

\begin{example}[Extracting the Site from Code]
\label{ex:site-extraction}
Consider:
\begin{lstlisting}
def safe_div(a: int, b: int) -> float:
    if b == 0:
        return float('inf')
    elif b > 0:
        return a / b
    else:
        return -(a / (-b))
\end{lstlisting}
The refinement site for this function has objects:
\begin{align*}
  \varphi_1 &\defeq (b = 0) \\
  \varphi_2 &\defeq (b > 0) \\
  \varphi_3 &\defeq (b < 0) \qquad [\text{implicit \texttt{else}}]
\end{align*}
The cover $\{\varphi_1, \varphi_2, \varphi_3\}$ is exhaustive:
$\varphi_1 \vee \varphi_2 \vee \varphi_3 \equiv (b = 0 \vee b > 0
\vee b < 0) \equiv \top$.  Z3 verifies this in microseconds.

The overlaps are:
$\varphi_1 \wedge \varphi_2 \equiv (b = 0 \wedge b > 0) \equiv \bot$
(disjoint), and similarly for all pairs.  No compatibility proofs are
needed.
\end{example}

\begin{example}[Non-Disjoint Cover: \texttt{bool} $\subset$ \texttt{int}]
\label{ex:bool-subtype}
Consider:
\begin{lstlisting}
def to_int(x):
    if isinstance(x, bool):
        return 1 if x else 0
    elif isinstance(x, int):
        return x
    elif isinstance(x, str):
        return len(x)
\end{lstlisting}
The cover $\{\varphi_\mathsf{bool}, \varphi_\mathsf{int},
\varphi_\mathsf{str}\}$ is \emph{not} disjoint:
$\varphi_\mathsf{bool} \wedge \varphi_\mathsf{int}$ is satisfiable
because every \texttt{bool} is an \texttt{int} in Python.
Z3 detects this via the subtyping morphism
$\mathsf{bool} \hookrightarrow \mathsf{int}$.

The compiler must then verify that the two branches \emph{agree} on
the overlap: for $x$ of type \texttt{bool}, the bool-branch returns
\texttt{1 if x else 0}, and the int-branch would return $x$ (which
is $1$ or $0$ for booleans).  The compatibility proof is:
$\forall x : \mathsf{bool}.\; (1\ \mathit{if}\ x\ \mathit{else}\ 0) = x$,
which Z3 verifies over the two-element domain $\{0, 1\}$.

Without the refinement site modeling $\mathsf{bool} \subseteq
\mathsf{int}$, this overlap would be silently missed, producing a
potentially unsound proof.
\end{example}


% ═══════════════════════════════════════════════════════════════
\section{What Cohomology Adds Beyond Cubical Type Theory}
\label{sec:what-cohomology-adds}

This section makes the central technical argument of this chapter.
We identify six concrete ways in which the cohomological machinery
provides capabilities that cubical type theory alone does not,
focusing on \emph{practical} gains in tractability and expressivity.


% ──────────────────────────────────────────────
\subsection{Tractability via Proof Decomposition}
\label{sec:decomposition}

The fundamental contribution of cohomology is
\textbf{proof decomposition}: given a global proof obligation
$\Gamma \vdash f \equiv_A g$, decompose it into local obligations
$\Gamma, x : \varphi_i \vdash \restr{f}{\varphi_i} \equiv_{A_i}
\restr{g}{\varphi_i}$ plus gluing conditions.

Cubical type theory provides \emph{paths} (equalities), \emph{transport}
(proof reuse), and \emph{function extensionality} (pointwise $\to$ global).
What it does \emph{not} provide is a mechanism to \emph{choose the
decomposition}.  In cubical type theory, a proof of $f = g$ is a single
path $p : \interval \to (A \to B)$ with $p(\mathsf{0}_\interval) = f$
and $p(\mathsf{1}_\interval) = g$.  The path must handle \emph{all}
inputs simultaneously---the interval does not decompose by input case.

Cohomological descent adds exactly this missing piece:

\begin{theorem}[Decomposition Theorem]
\label{thm:decomposition}
Let $f, g : A \to B$ and $\{\varphi_i\}_{i=1}^n$ be a cover of
$A$ (i.e., $\bigvee_i \varphi_i = \top$).  If:
\begin{enumerate}
\item For each $i$, there is a proof $p_i : \Path_{A_i \to B}
  (\restr{f}{\varphi_i}, \restr{g}{\varphi_i})$
  (local proof on fiber $\varphi_i$);
\item For each $i, j$ with $\varphi_i \wedge \varphi_j$ satisfiable,
  $\restr{p_i}{\varphi_i \wedge \varphi_j} =
  \restr{p_j}{\varphi_i \wedge \varphi_j}$
  (overlap compatibility);
\item $\Hone(\covers, \presheaf_{\Path}) = 0$ (cohomological obstruction
  vanishes);
\end{enumerate}
then there exists a global proof $p : \Path_{A \to B}(f, g)$.
\end{theorem}

\begin{proof}
By the descent theorem (Theorem~\ref{thm:c4-descent-full}), the
presheaf $\presheaf_\Path$ satisfies the sheaf condition.  The
compatible family $(p_i)$ with cocycle data from condition~(2)
defines a descent datum.  Condition~(3) ensures the cocycle is
a coboundary, so the descent datum extends to a global section.
\end{proof}

The key insight is condition~(3): for \emph{propositional} path types
(i.e., types where any two proofs of the same equality are themselves
equal---which holds for all path types in HoTT by Hedberg's theorem
when the type $B$ has decidable equality), $\Hone$ is \emph{automatically
trivial}.

\begin{proposition}[$\Hone = 0$ for Path Types over Decidable Types]
\label{prop:h1-trivial}
If $B$ has decidable equality (which includes $\Nat$, $\Int$,
$\mathsf{String}$, $\Bool$, and all finite types), then
$\Hone(\covers, \presheaf_{\Path_B}) = 0$ for any cover $\covers$.
\end{proposition}

\begin{proof}
The path type $\Path_B(a, b)$ is propositional when $B$ has decidable
equality (Hedberg's theorem).  A propositional presheaf over a
finite site has trivial $\Hone$ because the only cocycle is the
constant zero cocycle, which is always a coboundary.
\end{proof}

\begin{remark}
For higher types (types where $\Path$ is not propositional, e.g.,
types with non-trivial $\pi_2$), $\Hone$ can be non-trivial.  In
the Python setting, all runtime types have decidable equality, so
Proposition~\ref{prop:h1-trivial} applies universally.  The $\Hone$
computation serves as a \emph{sanity check}, not a real obstruction.
However, it would become substantive if $\mathrm{C}^4$ were extended
to verify programs manipulating higher-order abstract types (e.g.,
functors, monads as types with non-trivial homotopy structure).
\end{remark}


% ──────────────────────────────────────────────
\subsection{Complexity Reduction: Exponential to Linear}
\label{sec:complexity}

The practical impact of Theorem~\ref{thm:decomposition} is a
\emph{complexity reduction} in Z3 solving time.

\begin{theorem}[Complexity of Cohomological Descent]
\label{thm:complexity}
Let $f$ be a Python function with $n$ branches guarded by
predicates $\varphi_1, \ldots, \varphi_n$, and let $S$ be a
specification expressible as a refinement predicate on outputs.
\begin{enumerate}
\item \textbf{Monolithic approach:} The single Z3 query
  $\forall \vec{x}.\; \bigwedge_{i=1}^n (\varphi_i(\vec{x})
  \Rightarrow S(\vec{x}, f_i(\vec{x})))$
  has $n$ universal clauses that may interact combinatorially.
  In the worst case, Z3's DPLL(T) explores $O(2^n)$
  case splits.

\item \textbf{Cohomological approach:} The $n + \binom{n}{2} + 1$
  queries are:
  \begin{itemize}
  \item $1$ exhaustiveness check:
    $\neg(\bigvee_i \varphi_i)$ is UNSAT.
  \item $n$ fiber checks:
    $\neg(\varphi_i \Rightarrow S(\vec{x}, f_i(\vec{x})))$ is UNSAT,
    one per fiber.
  \item $\leq \binom{n}{2}$ overlap checks:
    $\neg(\varphi_i \wedge \varphi_j \Rightarrow
    f_i(\vec{x}) = f_j(\vec{x}))$ is UNSAT,
    one per non-disjoint pair (skipped for disjoint pairs).
  \end{itemize}
  Each individual query involves only one or two branches, so Z3
  handles it in $O(1)$ time (no inter-branch case splits).
  Total: $O(n^2)$ simple queries.

\item \textbf{Disjoint cover optimization:} When all
  $\varphi_i \wedge \varphi_j \equiv \bot$ (which holds for
  typical \texttt{if}/\texttt{elif}/\texttt{else} chains),
  the overlap checks vanish.  Total: $O(n)$ queries.
\end{enumerate}
\end{theorem}

\begin{proof}
Part (1): the monolithic formula has $n$ implications under a
universal quantifier.  Z3's case-splitting procedure may branch
on which $\varphi_i$ holds, yielding $2^n$ branches in the
worst case (though heuristics often do better).

Part (2): each fiber check $\varphi_i \Rightarrow S(\vec{x},
f_i(\vec{x}))$ involves exactly one branch function $f_i$ under
exactly one predicate $\varphi_i$.  There is no interaction with
other branches.  Similarly, each overlap check involves exactly
two branch functions under the conjunction of two predicates.

Part (3): if all pairs are disjoint, $\binom{n}{2}$ disjointness
checks (each trivially UNSAT by Z3) replace the overlap checks.
The disjointness checks are $O(1)$ each.  The total becomes
$1 + n + 0 = O(n)$ substantive queries plus $\binom{n}{2}$
trivial ones.
\end{proof}

\begin{example}[Concrete Speedup]
\label{ex:speedup}
Consider a type-dispatch function with $n = 10$ \texttt{isinstance}
branches (not uncommon in parsers, serializers, ORMs):
\begin{itemize}
\item Monolithic: one Z3 query with 10 universally quantified
  implications.  Z3 may explore up to $2^{10} = 1024$ case splits.
\item Cohomological: 10 fiber checks + 1 exhaustiveness check.
  Each fiber check is a single-branch implication—Z3 solves
  in microseconds.  Total wall time: $\sim 10 \times 50\mu s
  = 500\mu s$ vs.\ potentially $> 100ms$ monolithic.
\end{itemize}
This is a \textbf{200×} speedup, and the gap grows exponentially
with $n$.
\end{example}


% ──────────────────────────────────────────────
\subsection{The Cover Algebra}
\label{sec:cover-algebra}

Cubical type theory has no native notion of ``cover'' or
``decomposition strategy.''  You have paths, and you compose them.
The choice of how to decompose a proof is external to the theory.

Cohomology internalizes this choice as a \emph{cover algebra}:

\begin{definition}[Cover Operations]
\label{def:cover-ops}
Given covers $\covers = \{\varphi_i\}$ and $\covers' = \{\psi_j\}$:
\begin{enumerate}
\item \textbf{Refinement:} $\covers' \preceq \covers$ (``$\covers'$
  refines $\covers$'') if every $\psi_j$ is contained in some
  $\varphi_i$.  A proof via $\covers$ can be \emph{refined} to
  a proof via $\covers'$.

\item \textbf{Product:} $\covers \times \covers' \defeq
  \{\varphi_i \wedge \psi_j\}_{i,j}$.  This arises when a
  function $f$ calls a function $g$ and both have
  \texttt{isinstance} dispatch.  The proof of $f$ must
  consider all combinations of $f$'s and $g$'s branches.

\item \textbf{Join:} $\covers \vee \covers'$ is the coarsest
  common refinement.

\item \textbf{Restriction:} $\restr{\covers}{\psi} \defeq
  \{\varphi_i \wedge \psi\}_i$, the cover restricted to a
  sub-domain.
\end{enumerate}
All operations are decidable over Z3-decidable predicates:
refinement is $\forall j.\ \exists i.\ \psi_j \Rightarrow \varphi_i$;
product is conjunction; restriction is conjunction.
\end{definition}

\begin{remark}
Cubical type theory's interval $\interval$ and face formulas
$\varphi$ provide a superficially similar algebra, but it operates
on \emph{proof dimensions} (``along which dimension of the cube
are we composing?''), not on \emph{input domains} (``which subset
of inputs are we proving for?'').  These are orthogonal:
\begin{itemize}
\item The cubical interval decomposes \emph{how} we prove
  (which proof step).
\item The refinement cover decomposes \emph{what} we prove
  (which inputs).
\end{itemize}
$\mathrm{C}^4$ uses \emph{both}: the interval parameterizes proof
steps; the cover parameterizes input fibers.  The combination is
strictly more powerful than either alone.
\end{remark}


% ──────────────────────────────────────────────
\subsection{$\Hone$ as Proof Obstruction}
\label{sec:h1-obstruction}

When a proof \emph{cannot} be decomposed locally, the cohomological
framework tells you \emph{why}:

\begin{theorem}[$\Hone$ Obstruction]
\label{thm:h1-obstruction}
Let $\covers = \{\varphi_i\}$ be a cover with non-trivial overlaps.
If $\Hone(\covers, \presheaf_\Path) \neq 0$, then:
\begin{enumerate}
\item No compatible family of local proofs extends to a global proof
  via this cover.
\item The obstruction lies in the \emph{overlap structure}, not in
  the individual fibers.
\item A \emph{finer} cover $\covers' \preceq \covers$ may have
  $\Hone(\covers', \presheaf_\Path) = 0$, enabling descent.
\end{enumerate}
\end{theorem}

This is information that cubical type theory simply does not
provide.  In cubical TT, if \textsf{hcomp} fails, you get no
diagnostic about \emph{why} the faces don't glue.  With
$\Hone$, you get a precise obstruction class telling you
which overlaps are problematic.

\begin{example}[Non-Trivial $\Hone$ in Practice]
\label{ex:nontrivial-h1}
While Proposition~\ref{prop:h1-trivial} shows $\Hone = 0$ for
equality proofs on decidable types, non-trivial $\Hone$ can
arise when verifying \emph{higher-order} properties.

Consider verifying that a function $f$ is \emph{monotone}:
$\forall x \leq y.\ f(x) \leq f(y)$.  Monotonicity is a
\emph{relation} between pairs of inputs, not a pointwise
property.  The presheaf $\presheaf_\mathsf{mono}$ assigns
to each fiber $\varphi$ the set of proofs that $f$ is monotone
on $\{x \mid \varphi(x)\}$.

Suppose the cover is $\{x \leq 0, x \geq 0\}$ with overlap
$\{x = 0\}$.  On each fiber, $f$ may be monotone (increasing on
negatives, increasing on positives).  But the overlap compatibility
requires that the \emph{slopes match at $x = 0$}: if $f$ is
decreasing-then-increasing (a ``V'' shape), the local monotonicity
proofs don't glue.  The obstruction is exactly $\Hone \neq 0$:
the cocycle on the overlap records the ``direction mismatch'' at
the seam.

In this case, the correct action is to \emph{refine the cover}
or \emph{change the property being proved} (e.g., prove piecewise
monotonicity rather than global monotonicity).  $\Hone$ directs
the proof strategy.
\end{example}


% ──────────────────────────────────────────────
\subsection{The Spectral Sequence as Proof Strategy}
\label{sec:spectral}

The \v{C}ech spectral sequence provides a \emph{systematic strategy}
for finding proofs.  This is a structured proof search that neither
cubical type theory nor classical Hoare logic provides.

\begin{definition}[\v{C}ech Spectral Sequence for Proofs]
The pages of the spectral sequence correspond to stages of the
proof:
\begin{enumerate}
\item[$E_0$:] \textbf{Local data.}
  Compute the behavior of $f$ on each fiber $\varphi_i$.
  This is function restriction: $\restr{f}{\varphi_i}$.

\item[$E_1$:] \textbf{Local proofs.}
  For each fiber $\varphi_i$, attempt to prove
  $\restr{f}{\varphi_i} \equiv \restr{g}{\varphi_i}$.
  These are the fiber VCs.
  $E_1^{p,0} = \Hzero(\covers, \presheaf)$ = the local sections.

\item[$E_2$:] \textbf{Compatibility.}
  Check overlap agreement.  Compute
  $E_2^{p,q} = \check{H}^p(E_1^{\bullet, q})$.
  If $E_2^{1,0} = \Hone = 0$, the local proofs glue.

\item[$E_\infty$:] \textbf{Global proof.}
  If the spectral sequence collapses at $E_2$ (which it does
  for propositional coefficients), the global proof exists.
\end{enumerate}
\end{definition}

\begin{remark}
This strategy is \emph{algorithmic}:
\begin{enumerate}
\item Extract the cover from the code's control flow (AST analysis).
\item Run Z3 on each fiber VC (parallelizable).
\item Run Z3 on each overlap VC (parallelizable).
\item Compute $\Hone$ from the overlap graph (linear-time BFS).
\item If all pass, emit the global proof by descent.
\end{enumerate}
No human intervention is needed for the common case of disjoint
\texttt{if}/\texttt{elif}/\texttt{else} chains.  For non-disjoint
covers, the overlap VCs may require human-provided compatibility
lemmas, but $\Hone$ computation tells you \emph{exactly which
overlaps} need attention.
\end{remark}


% ──────────────────────────────────────────────
\subsection{Incremental Verification}
\label{sec:incremental}

When a single branch of a function changes, the cohomological
framework enables \emph{incremental reverification}:

\begin{proposition}[Incremental Reverification]
\label{prop:incremental}
Let $f$ be verified via cohomological descent over cover
$\{\varphi_i\}$.  If branch $k$ is modified (producing $f'$
with $\restr{f'}{\varphi_k} \neq \restr{f}{\varphi_k}$ but
$\restr{f'}{\varphi_i} = \restr{f}{\varphi_i}$ for $i \neq k$),
then:
\begin{enumerate}
\item Only the fiber VC for $\varphi_k$ must be re-checked.
\item Only the overlap VCs involving $\varphi_k$ must be re-checked.
\item All other fiber and overlap VCs remain valid.
\end{enumerate}
\end{proposition}

\begin{proof}
The descent proof decomposes as a product of local proofs.
Modifying one factor ($p_k$) does not affect the others
($p_i$ for $i \neq k$).  The overlap proofs
$p_{ij}$ for $i,j \neq k$ are unchanged.
\end{proof}

In cubical type theory without cohomological decomposition, there
is no analog: the global path $p : \interval \to (A \to B)$ is a
\emph{single} object, and modifying one branch requires re-proving
the entire path.


% ──────────────────────────────────────────────
\subsection{Composition of Covers: The K\"unneth Formula}
\label{sec:kunneth}

When function $f$ calls function $g$, and both have
\texttt{isinstance} dispatch, the combined proof must handle
all combinations of their branches.

\begin{proposition}[Compositional Descent]
\label{prop:compositional}
Let $f$ have cover $\covers_f = \{\varphi_i\}_{i=1}^m$ and
$g$ have cover $\covers_g = \{\psi_j\}_{j=1}^n$.  The
composed function $f \circ g$ (or more generally, any expression
mixing $f$ and $g$) has the product cover
$\covers_f \times \covers_g = \{\varphi_i \wedge \psi_j\}$.

The proof of the composed expression decomposes into $mn$ fiber
VCs (one per product fiber), plus overlap checks.  The K\"unneth
formula gives:
\[
  \Hone(\covers_f \times \covers_g, \presheaf) \;\cong\;
  (\Hone(\covers_f) \otimes \Hzero(\covers_g)) \oplus
  (\Hzero(\covers_f) \otimes \Hone(\covers_g))
\]
so if \emph{both} individual covers have $\Hone = 0$, the
product cover also has $\Hone = 0$.
\end{proposition}

This compositionality is essential for scaling verification
beyond single functions.  Cubical type theory can compose
paths ($p \cdot q$), but has no mechanism for composing
\emph{decomposition strategies}.


% ═══════════════════════════════════════════════════════════════
\section{The Cubical Structure of Python Proofs}
\label{sec:cubical-structure}

We now explain how the cubical structure of $\mathrm{C}^4$ interacts
with the cohomological decomposition to produce a verification
framework more powerful than either alone.


\subsection{Proofs as Cubical Paths}
\label{sec:proofs-paths}

Every proof term $p : f \equiv g$ is a \textbf{cubical path}:
a term $p : \interval \to (A \to B)$ with $p(\mathsf{0}_\interval) = f$
and $p(\mathsf{1}_\interval) = g$.  The \emph{face maps} extract
endpoints:
\[
  \partial^0(p) = p(\mathsf{0}_\interval) = f, \qquad
  \partial^1(p) = p(\mathsf{1}_\interval) = g.
\]

In the VC compiler, the face-map check is a fundamental invariant:
every proof must have correct faces.  For \textsf{Trans}$(p, q)$,
the \textbf{seam condition} is $\partial^1(p) = \partial^0(q)$:
the target of $p$ must equal the source of $q$, definitionally.


\subsection{The Unification: Descent = HComp}
\label{sec:descent-hcomp}

The central technical insight of the compiler is:

\begin{theorem}[Descent--HComp Equivalence]
\label{thm:descent-hcomp}
For a cover $\covers = \{\varphi_i\}_{i=1}^n$ with local proofs
$(p_i)$ and overlap proofs $(p_{ij})$:
\[
  \mathsf{desc}(\{p_i\}, \{p_{ij}\})
  \;=\;
  \hfilling^{i_1, \ldots, i_n}\,
  \Path_{A \to B}\,
  [\vec{\varphi} \mapsto \vec{p}]\,
  p_{\mathrm{base}}
\]
That is, the descent construction is a special case of
cubical homogeneous composition (\textsf{hcomp}) where:
\begin{itemize}
\item Each fiber $\varphi_k$ corresponds to a \emph{cubical face}
  $[i_k = 0]$ or $[i_k = 1]$ in a multi-dimensional cube.
\item The local proof $p_k$ is the partial boundary data on
  face $k$.
\item The overlap compatibility $p_{ij}$ is the \emph{edge agreement}
  between adjacent faces.
\item The Kan filling condition (``\textsf{hcomp} exists'')
  is exactly $\Hone = 0$.
\end{itemize}
\end{theorem}

\begin{proof}[Proof sketch]
Both constructions solve the same universal problem: given
compatible partial data on the boundary of a ``space,'' produce
a global section.  For descent, the space is the site; for
\textsf{hcomp}, it is the cube.  The equivalence follows from
the fact that the Čech nerve of a finite cover is (the geometric
realization of) a simplicial set that embeds into a cubical set,
and the Kan condition for the cubical set implies the sheaf
condition for the site.
\end{proof}

\begin{remark}
This unification is not just theoretical elegance; it has a
practical consequence: the VC compiler uses a \textbf{single
implementation} for both descent and \textsf{hcomp} proof terms.
The function \texttt{\_compile\_cubical\_filling} handles both,
differing only in how the ``faces'' and ``edges'' are specified
(refinement predicates for descent; cubical face names for
\textsf{hcomp}).
\end{remark}


\subsection{Transport as Proof Reuse}
\label{sec:transport-reuse}

The cubical \textsf{transp} operation moves proofs along type paths:
\[
  \frac{\Gamma \vdash p : \Path_\Type(A, B) \qquad
        \Gamma \vdash t : A}
       {\Gamma \vdash \transp^i\, p\, t : B}
\]

For Python verification, this has a powerful practical interpretation:
\emph{if you proved a property for one implementation and have a path
to another implementation, you get the property for free}.

\begin{example}[Transport Between Algorithms]
\label{ex:transport-algo}
Suppose we proved:
\begin{itemize}
\item $p_\mathsf{merge} :$ \texttt{merge\_sort} satisfies
  ``returns sorted permutation.''
\item $q :$ \texttt{merge\_sort} $\equiv$ \texttt{quick\_sort}
  (via the deterministic specification uniqueness theorem,
  Theorem~\ref{thm:det-spec}).
\end{itemize}
Then $\transp^i\, q\, p_\mathsf{merge}$ gives us:
\texttt{quick\_sort} satisfies ``returns sorted permutation''
--- without re-proving from scratch.

The VC for this transport is:
\begin{enumerate}
\item $q$ is a valid path (both algorithms satisfy the spec).
\item $p_\mathsf{merge}$ is a valid proof.
\item The transported proof has correct faces:
  $\partial^0 = $ ``\texttt{merge\_sort} satisfies spec'' (from
  $p_\mathsf{merge}$) and $\partial^1 = $
  ``\texttt{quick\_sort} satisfies spec'' (the transported result).
\end{enumerate}
\end{example}

Transport is orthogonal to cohomological descent: descent decomposes
proofs \emph{across fibers}; transport moves proofs
\emph{across types/implementations}.  Using both:
\begin{enumerate}
\item Prove property $P$ for \texttt{merge\_sort} via descent
  (decomposing over input fibers: empty list, singleton, general).
\item Transport to \texttt{quick\_sort} via a shared specification
  path.
\end{enumerate}
Total cost: $n$ fiber VCs for \texttt{merge\_sort} +
1 specification proof for \texttt{quick\_sort} +
1 transport.  Without these tools: a fresh proof from scratch
for each algorithm.


\subsection{HComp for Concurrent Changes}
\label{sec:hcomp-concurrent}

The \textsf{hcomp} (homogeneous composition) operation fills a
\emph{cube} from compatible face data.  In practice, this handles
the case where multiple proof paths exist between the same endpoints.

\begin{example}[Refactoring Consistency]
\label{ex:refactoring}
Suppose old code $f$ is refactored to new code $g$ via two
independent changes:
\begin{itemize}
\item Path $p$: change the algorithm (e.g., bubble sort $\to$
  merge sort).
\item Path $q$: change the data representation (e.g., list $\to$
  numpy array).
\end{itemize}
These produce a \emph{square}:
\[
\begin{tikzcd}
  f_{\mathrm{list}} \ar[r, "p_{\mathrm{list}}"]
    \ar[d, "q_{\mathrm{bubble}}"']
    & g_{\mathrm{list}} \ar[d, "q_{\mathrm{merge}}"] \\
  f_{\mathrm{array}} \ar[r, "p_{\mathrm{array}}"']
    & g_{\mathrm{array}}
\end{tikzcd}
\]
\textsf{hcomp} verifies that this square \emph{commutes}: applying
the changes in either order gives the same result.  The VCs are:
\begin{enumerate}
\item $p_\mathrm{list}$ is valid (algorithm change on lists).
\item $p_\mathrm{array}$ is valid (algorithm change on arrays).
\item $q_\mathrm{bubble}$ is valid (representation change for
  bubble sort).
\item $q_\mathrm{merge}$ is valid (representation change for
  merge sort).
\item Corners match: $\partial^1(p_\mathrm{list}) =
  \partial^0(q_\mathrm{merge})$, etc.
\end{enumerate}
\end{example}


% ═══════════════════════════════════════════════════════════════
\section{The Verification Condition Compiler}
\label{sec:vc-compiler}

We now give the formal compilation rules.  Each proof term maps to
a finite set of VCs.  The compiler is implemented as
\texttt{c4\_compiler.py}.


\subsection{The Compilation Judgment}

\begin{definition}[VC Compilation]
\label{def:vc-compilation}
The \textbf{compilation judgment} has the form:
\[
  \Gamma \vdash p : f \equiv_A g
  \;\leadsto\;
  \{\mathsf{VC}_1, \ldots, \mathsf{VC}_k\}, \;\; \mathsf{Trust}
\]
where:
\begin{itemize}
\item $\Gamma$ is the typing context (variable declarations).
\item $p$ is a proof term from the $\mathrm{C}^4$ proof language.
\item $f, g$ are OTerms (the LHS and RHS of the equality).
\item $\mathsf{VC}_i$ is a Z3 formula plus metadata (source rule,
  status, trust).
\item $\mathsf{Trust}$ is a \emph{provenance set} recording which
  external sources the proof depends on.
\end{itemize}
\end{definition}

\begin{definition}[Soundness of VC Compilation]
\label{def:vc-soundness}
The compiler is \textbf{sound} if: whenever all VCs are valid
(either Z3-verified or explicitly assumed), the proof term $p$
is a valid $\mathrm{C}^4$ judgment.  Formally:
\[
  \left(\forall i.\; \mathsf{VC}_i \text{ is valid}\right)
  \;\Longrightarrow\;
  \Gamma \vdash p : f \equiv_A g \text{ in } \mathrm{C}^4.
\]
\end{definition}


\subsection{Compilation Rules}

We give the compilation rule for each group of proof terms.

\paragraph{Kernel rules (Trust = KERNEL).}

\begin{align}
  \mathsf{Refl} &\;\leadsto\;
    \{\mathsf{VC}(\text{``}\mathit{norm}(f) \stackrel{?}{=}
    \mathit{norm}(g)\text{''})\}
    \tag{no Z3; syntactic comparison after normalization}
    \\[4pt]
  \mathsf{Sym}(p) &\;\leadsto\;
    \mathsf{compile}(p, g, f)
    \tag{swap LHS/RHS, compile inner}
    \\[4pt]
  \mathsf{Trans}(p, q) &\;\leadsto\;
    \mathsf{compile}(p) \cup \mathsf{compile}(q) \cup
    \{\mathsf{VC}(\partial^1(p) \equiv \partial^0(q))\}
    \tag{seam check}
    \\[4pt]
  \mathsf{Cong}(h, p) &\;\leadsto\;
    \mathsf{compile}(p) \cup
    \{\mathsf{VC}(\text{``congruence under } h\text{''})\}
\end{align}

\paragraph{Descent rules (Trust $\ni$ Z3).}

The unified cubical filling compilation:
\begin{align}
  &\mathsf{RefinementDescent}(\{(\varphi_i, p_i)\}, \{p_{ij}\})
  \;\leadsto\;
  \notag \\
  &\quad \underbrace{\{\mathsf{VC}(\bigvee_i \varphi_i)\}}_\text{exhaustive}
  \;\cup\;
  \underbrace{\bigcup_i \mathsf{compile}(p_i)}_\text{fiber VCs}
  \;\cup\;
  \underbrace{\bigcup_{i < j,\; \varphi_i \wedge \varphi_j \neq \bot}
    \mathsf{compile}(p_{ij})}_\text{overlap VCs}
  \;\cup\;
  \underbrace{\{\mathsf{VC}(\Hone = 0)\}}_\text{Kan condition}
\end{align}
The same compilation is used for \textsf{GluePath}, \textsf{CechGluing},
\textsf{FiberDecomposition}, and \textsf{HComp} — differing only in how
the faces and overlaps are specified.

\paragraph{Z3 rules (Trust = Z3).}

\begin{align}
  \mathsf{Z3Discharge}(\varphi, \sigma) &\;\leadsto\;
    \{\mathsf{VC}(\varphi,\ \text{check } \neg\varphi \text{ is UNSAT})\}
\end{align}

\paragraph{Induction rules (Trust $\ni$ Z3).}

\begin{align}
  &\mathsf{NatInduction}(p_0, p_S)
  \;\leadsto\;
  \underbrace{\mathsf{compile}(p_0)}_\text{base}
  \;\cup\;
  \underbrace{\mathsf{compile}(p_S)}_\text{step}
  \;\cup\;
  \{\mathsf{VC}(\text{``well-founded on } \Nat\text{''})\}
\end{align}

\paragraph{Import rules (Trust $\ni$ LEAN or LIBRARY).}

\begin{align}
  \mathsf{MathLibAxiom}(\mathit{name}) &\;\leadsto\;
    \{\mathsf{VC}(\text{``assumed: } \mathit{name}\text{''}, \mathsf{LEAN})\}
  \\[4pt]
  \mathsf{LibraryAxiom}(\mathit{lib}, \mathit{ax}) &\;\leadsto\;
    \{\mathsf{VC}(\text{``assumed: } \mathit{lib}.\mathit{ax}\text{''},
    \mathsf{LIBRARY})\}
\end{align}

\paragraph{Transport rules.}

\begin{align}
  &\mathsf{Transport}(q, p_\mathsf{src}) \;\leadsto\;
  \mathsf{compile}(q) \;\cup\;
  \mathsf{compile}(p_\mathsf{src}) \;\cup\;
  \notag\\
  &\qquad \{\mathsf{VC}(\text{``target predicate } \Rightarrow
  \text{ source predicate''}, \text{Z3})\}
\end{align}


% ═══════════════════════════════════════════════════════════════
\section{Trust Provenance}
\label{sec:trust}

Classical verification systems assign a single ``trust level'' to
each proof: verified, assumed, or unknown.  $\mathrm{C}^4$
tracks \emph{provenance}:

\begin{definition}[Trust Provenance]
\label{def:trust-provenance}
A \textbf{trust provenance} is a pair $(\mathcal{S}, \mathcal{A})$
where:
\begin{itemize}
\item $\mathcal{S} \subseteq \{\mathsf{KERNEL}, \mathsf{Z3},
  \mathsf{LEAN}, \mathsf{LIBRARY}, \mathsf{AXIOM},
  \mathsf{ASSUMED}\}$ is the set of trust sources used.
\item $\mathcal{A}$ is a list of named assumptions (e.g.,
  ``\texttt{lean:Nat.add\_comm}'',
  ``\texttt{lib:sympy:expand\_add}'').
\end{itemize}
Composition of proofs composes provenance:
$(\mathcal{S}_1, \mathcal{A}_1) \otimes (\mathcal{S}_2, \mathcal{A}_2)
\defeq (\mathcal{S}_1 \cup \mathcal{S}_2,
\mathcal{A}_1 \cup \mathcal{A}_2)$.
\end{definition}

\begin{remark}[Why Not a Total Order?]
F* assigns trust as $\mathsf{Tot} > \mathsf{ML} > \mathsf{ALL}$
(a total order on effects).  Lean assigns trust as
$\mathsf{kernel} > \mathsf{tactic} > \mathsf{sorry}$.  Both are
total orders because each system has a single trust hierarchy.

$\mathrm{C}^4$ has \emph{multiple independent} trust sources:
a proof might use both a Lean theorem (machine-checked in a
different system) and a library axiom (human-asserted about
Python).  These are incomparable: neither ``dominates'' the other.
Making trust a set rather than an order lets the \emph{consumer}
decide what they trust, rather than the \emph{prover}.
\end{remark}


% ═══════════════════════════════════════════════════════════════
\section{F*-Style Annotation Binding}
\label{sec:binding}

The final component ensuring soundness of the compiler is
\textbf{annotation binding}: the requirement that a proof
annotation must \emph{structurally match} the code it claims
to describe.

\begin{definition}[Binding Judgment]
\label{def:binding}
An annotation $\mathcal{A}$ \textbf{binds} to source code $c$ if:
\begin{enumerate}
\item \textbf{Name match:} $\mathcal{A}$ names a function $f$;
  $c$ defines a function named $f$.
\item \textbf{Parameter match:} The parameter names and order in
  $\mathcal{A}$ equal those in $c$.
\item \textbf{Branch coverage:} Each \texttt{if}/\texttt{elif}
  /\texttt{else} branch in $c$ corresponds to a fiber in
  $\mathcal{A}$'s refinement cover.
\item \textbf{Source hash:} A hash of $c$ is stored in $\mathcal{A}$
  to detect staleness (optional but recommended).
\end{enumerate}
If binding fails, the proof is \textbf{rejected} regardless of
Z3 results.
\end{definition}

This mirrors F*'s mechanism where the type of a function IS
its specification, and the function body must inhabit that type.
In $\mathrm{C}^4$, the annotation IS the proof, and the code
must be the witness.


% ═══════════════════════════════════════════════════════════════
\section{Worked Examples}
\label{sec:worked-examples}


\subsection{Example 1: isinstance Dispatch}
\label{sec:ex-isinstance}

\begin{lstlisting}
def describe(x):
    if isinstance(x, int):
        return f"integer {x}"
    elif isinstance(x, str):
        return f"string of length {len(x)}"
    elif isinstance(x, list):
        return f"list with {len(x)} elements"
    else:
        return "unknown"
\end{lstlisting}

\textbf{Specification:} $\forall x.\; \mathit{describe}(x)$
starts with the correct type name.

\textbf{Refinement cover:}
$\covers = \{\varphi_\mathsf{int}, \varphi_\mathsf{str},
\varphi_\mathsf{list}, \varphi_\mathsf{else}\}$.

\textbf{VCs generated by the compiler:}
\begin{enumerate}
\item \emph{Exhaustiveness:}
  $\mathit{isinstance}(x, \mathsf{int}) \vee
  \mathit{isinstance}(x, \mathsf{str}) \vee
  \mathit{isinstance}(x, \mathsf{list}) \vee
  \varphi_\mathsf{else}$
  where $\varphi_\mathsf{else} = \neg(\text{int} \vee \text{str}
  \vee \text{list})$.
  Z3: valid (by construction of the else branch).

\item \emph{Fiber VCs:} 4 checks, one per branch.  Under the
  hypothesis $\varphi_i$, the return value starts with the
  correct type name.  Each is a simple string-prefix check
  (structural, not Z3).

\item \emph{Overlaps:} All pairs are disjoint (distinct
  \texttt{isinstance} types, assuming no subtype relationships
  among int, str, list).  No compatibility proofs needed.

\item \emph{$\Hone$:} 0 (disjoint cover, single component).
\end{enumerate}

\textbf{Trust:} KERNEL (structural checks only, no Z3 or
library axioms needed for this simple case).

\textbf{Comparison without cohomology:}
A monolithic proof would need to handle all 4 branches in a single
path term, threading the interval through the case analysis.  With
4 branches, this is manageable; with 20 branches (common in
serializers), it becomes prohibitive.


\subsection{Example 2: Library Axiom Transport}
\label{sec:ex-sympy}

\begin{lstlisting}
import sympy

def simplify_twice(expr):
    return sympy.simplify(sympy.simplify(expr))
\end{lstlisting}

\textbf{Specification:}
$\mathit{simplify\_twice}(e) = \mathit{sympy.simplify}(e)$.

\textbf{Proof term:}
\begin{lstlisting}[language={}]
LibraryAxiom(
  library="sympy",
  axiom_name="simplify_idempotent",
  instantiation={"e": OVar("expr")}
)
\end{lstlisting}

\textbf{VCs:} None (the library axiom is an assumption).

\textbf{Trust:} $\{\mathsf{LIBRARY}\}$, assumption:
``\texttt{lib:sympy:simplify\_idempotent}''.

The consumer of this proof sees exactly what was assumed:
sympy's \texttt{simplify} is idempotent.  If this assumption
is wrong (e.g., due to a sympy bug), the proof is invalid,
and the assumption list tells you precisely which assumption
failed.


\subsection{Example 3: Transport + Descent Combined}
\label{sec:ex-transport-descent}

\begin{lstlisting}
def abs_value(x: int) -> int:
    if x >= 0:
        return x
    else:
        return -x
\end{lstlisting}

\textbf{Specification:} $\mathit{abs\_value}(x) \geq 0$
for all $x : \Int$.

\textbf{Proof by cohomological descent:}
\[
  \mathsf{RefinementDescent}\left(
  \begin{array}{l}
    \text{cover: } \{\varphi_+ \defeq x \geq 0,\;
    \varphi_- \defeq x < 0\} \\
    \text{fiber } \varphi_+:\;
    \mathsf{Z3Discharge}(x \geq 0 \Rightarrow x \geq 0) \\
    \text{fiber } \varphi_-:\;
    \mathsf{Z3Discharge}(x < 0 \Rightarrow -x \geq 0)
  \end{array}\right)
\]

\textbf{VCs:}
\begin{enumerate}
\item Exhaustiveness: $(x \geq 0) \vee (x < 0)$.
  Z3: valid.  [\textsf{Z3}, 0.1ms]
\item Fiber $\varphi_+$: $x \geq 0 \Rightarrow x \geq 0$.
  Z3: valid (trivial).  [\textsf{Z3}, 0.05ms]
\item Fiber $\varphi_-$: $x < 0 \Rightarrow -x \geq 0$.
  Z3: valid ($x < 0 \Rightarrow -x > 0$).  [\textsf{Z3}, 0.1ms]
\item Overlap: $\varphi_+ \wedge \varphi_- \equiv (x \geq 0 \wedge x < 0)
  \equiv \bot$.  Disjoint.  [\textsf{Z3}, 0.05ms]
\item $\Hone = 0$ (disjoint, connected).  [\textsf{KERNEL}]
\end{enumerate}

\textbf{Trust:} $\{\mathsf{KERNEL}, \mathsf{Z3}\}$.
No assumptions.  Fully machine-checked.

\textbf{Now transport:} Suppose we also have:
\begin{lstlisting}
def abs_via_max(x: int) -> int:
    return max(x, -x)
\end{lstlisting}

The path $\mathit{abs\_value} \equiv \mathit{abs\_via\_max}$ can
be proved by descent (both return $x$ when $x \geq 0$ and $-x$
when $x < 0$, and $\max(x, -x)$ equals $x$ when $x \geq 0$ and
$-x$ when $x < 0$).

Transport gives: $\mathit{abs\_via\_max}(x) \geq 0$ for free,
reusing the proof for \texttt{abs\_value}.


\subsection{Example 4: Mathlib Path in a Python Proof}
\label{sec:ex-mathlib}

\begin{lstlisting}
def sum_first_n(n: int) -> int:
    total = 0
    for i in range(1, n + 1):
        total += i
    return total
\end{lstlisting}

\textbf{Specification:}
$\mathit{sum\_first\_n}(n) = n(n+1)/2$.

\textbf{Proof term:}
\[
  \mathsf{Trans}\left(
    \underbrace{\mathsf{LoopInvariant}(\ldots)}_
      {\text{total} = \sum_{i=1}^k i \text{ after } k \text{ iterations}},\;
    \underbrace{\mathsf{MathLibAxiom}(\texttt{Finset.sum\_range\_id\_eq\_sum})}_
      {\sum_{i=1}^n i = n(n+1)/2}
  \right)
\]

\textbf{VCs:}
\begin{enumerate}
\item Loop invariant init: $\mathit{total}_0 = 0 = \sum_{i=1}^0 i$.
  Z3: valid.
\item Loop invariant preservation:
  $\mathit{total}_k = \sum_{i=1}^k i \Rightarrow
  \mathit{total}_{k+1} = \sum_{i=1}^{k+1} i$.
  Z3: valid (given $\mathit{total}_{k+1} = \mathit{total}_k + (k+1)$).
\item Loop invariant exit:
  $\mathit{total}_n = \sum_{i=1}^n i$.
  Structural (from the invariant at termination).
\item Seam: $\partial^1(\text{loop proof}) = \partial^0(\text{Mathlib})$.
  Both sides: $\sum_{i=1}^n i$.  Definitional equality.
\item Mathlib axiom: \texttt{Finset.sum\_range\_id\_eq\_sum}.
  Assumed with \textsf{LEAN} trust.
\end{enumerate}

\textbf{Trust:} $\{\mathsf{Z3}, \mathsf{LEAN}\}$, assumption:
``\texttt{lean:Finset.sum\_range\_id\_eq\_sum}''.

This example shows the two systems working together: Z3 handles
the loop invariant (an arithmetic property), and Lean provides
the closed-form identity (a number-theoretic theorem).  Neither
alone would suffice.


\subsection{Example 5: Higher Coherence via HComp}
\label{sec:ex-hcomp}

Consider verifying that two independently developed patches
to a function are \emph{compatible}:
\begin{itemize}
\item Patch A: adds input validation.
\item Patch B: optimizes the hot path.
\end{itemize}

The \textsf{HComp} proof fills the square:
\[
\begin{tikzcd}[column sep=4em, row sep=3em]
  f_\text{original}
    \ar[r, "p_A\text{: add validation}"]
    \ar[d, "p_B\text{: optimize}"']
  & f_{A}
    \ar[d, "p_B'\text{: optimize (with validation)}"] \\
  f_{B}
    \ar[r, "p_A'\text{: add validation (optimized)}"']
  & f_{AB}
\end{tikzcd}
\]

\textbf{VCs from HComp:}
\begin{enumerate}
\item All four edge proofs are valid.
\item Corners match: $\partial^1(p_A) = \partial^0(p_B')$,
  $\partial^1(p_B) = \partial^0(p_A')$, etc.
\item The filled square commutes:
  $p_B' \circ p_A = p_A' \circ p_B$ (applying both patches in
  either order gives the same result).
\end{enumerate}

This is a genuinely \emph{2-dimensional} proof that cubical
type theory handles natively.  Without cubical structure,
verifying patch compatibility would require an ad hoc
argument.  Without cohomological structure, verifying each
edge proof would require monolithic reasoning.  Together,
they give a clean compositional verification of concurrent
changes.


% ═══════════════════════════════════════════════════════════════
\section{Scalability}
\label{sec:scalability}

We analyze the scalability of the cohomological approach on
real Python codebases.

\begin{theorem}[Linear Scaling for Disjoint Dispatch]
\label{thm:linear-scaling}
Let $L$ be a Python library with $m$ functions, each having
at most $n$ disjoint \texttt{isinstance} branches, and let
each branch body be Z3-decidable.  Then:
\begin{enumerate}
\item Total VCs: $\leq m \cdot (n + 1)$ (one exhaustiveness
  + $n$ fiber checks per function).
\item Total Z3 time: $O(m \cdot n \cdot t_\mathsf{Z3})$ where
  $t_\mathsf{Z3}$ is the per-query time (typically $< 1$ms
  for QF\_LIA).
\item Total compile time: $O(m \cdot n)$ (dominated by Z3).
\end{enumerate}
For a library with $m = 500$ functions and $n = 5$ branches
each, this is $\sim 2500$ Z3 queries, completing in
$< 1$ second.
\end{theorem}

\begin{remark}[Where Scaling Breaks Down]
Scaling is \emph{not} linear when:
\begin{enumerate}
\item \textbf{Non-disjoint covers:} Overlap checks add $O(n^2)$
  per function.  For $n = 20$ with all pairwise overlaps,
  this is 190 overlap checks per function.
\item \textbf{Non-Z3-decidable fibers:} When a fiber involves
  library operations (e.g., ``$\mathit{sympy.is\_polynomial}(e)$''),
  the fiber VC must be discharged by library axioms or oracle,
  not Z3.
\item \textbf{Higher-order functions:} Functions taking function
  arguments create covers over \emph{function spaces}, which
  are not directly Z3-decidable.
\item \textbf{Recursive functions:} Require induction, which
  adds structural complexity beyond simple descent.
\end{enumerate}
These limitations are fundamental: they arise from the
undecidability of the halting problem, not from the proof
framework.  Cohomological descent maximizes the \emph{decidable
fraction}; the rest requires human-provided lemmas, library
axioms, or oracle judgment.
\end{remark}


% ═══════════════════════════════════════════════════════════════
\section{Comparison with Existing Systems}
\label{sec:comparison}

\begin{center}
\begin{longtable}{@{}lcccc@{}}
\toprule
\textbf{Feature}
  & \textbf{F*}
  & \textbf{Lean 4}
  & \textbf{Dafny}
  & $\mathbf{C^4}$ \\
\midrule
Dependent types & \checkmark & \checkmark & partial & \checkmark \\
Cubical paths & --- & --- & --- & \checkmark \\
Transport/univalence & --- & --- & --- & \checkmark \\
Cohomological descent & --- & --- & --- & \checkmark \\
Refinement types & \checkmark & partial & \checkmark & \checkmark \\
Z3 integration & \checkmark & --- & \checkmark & \checkmark \\
Mathlib access & --- & \checkmark & --- & \checkmark \\
Library axioms & partial & partial & partial & \checkmark \\
Python target & --- & --- & --- & \checkmark \\
Proof decomposition & manual & manual & manual & automatic \\
Incremental reverify & --- & partial & --- & \checkmark \\
Trust provenance & total order & total order & binary & set \\
\bottomrule
\end{longtable}
\end{center}

\paragraph{Key differentiators of $\mathrm{C}^4$:}

\begin{enumerate}
\item \textbf{Automatic proof decomposition.}
  F*, Lean, and Dafny require the user to manually structure
  proofs.  $\mathrm{C}^4$ extracts the decomposition from the
  code's control flow and applies descent automatically.

\item \textbf{Cubical + cohomological.}
  No existing system combines cubical paths with sheaf-theoretic
  descent.  Each alone is useful; together they handle both
  proof reuse (transport) and proof decomposition (descent).

\item \textbf{Trust provenance as a set.}
  Other systems use binary (verified/not) or total-order trust.
  $\mathrm{C}^4$ tracks exactly which assumptions a proof depends
  on, letting consumers make informed trust decisions.

\item \textbf{Mathlib + Z3 in one system.}
  Lean has Mathlib but no Z3.  F* and Dafny have Z3 but no
  Mathlib.  $\mathrm{C}^4$ uses both: Z3 for decidable fragments,
  Mathlib for mathematical theorems, library axioms for
  Python-specific properties.
\end{enumerate}


% ═══════════════════════════════════════════════════════════════
\section{Conclusion}
\label{sec:compiler-conclusion}

This chapter demonstrated that the cohomological component of
$\mathrm{C}^4$ provides three concrete gains beyond what cubical
type theory alone offers:

\begin{enumerate}
\item \textbf{Tractability:} Cohomological descent reduces
  proof complexity from exponential to linear for branching code,
  with each local obligation in a Z3-decidable fragment.

\item \textbf{Expressivity:} The cover algebra, $\Hone$ obstruction,
  and spectral sequence provide proof-search structure that cubical
  type theory does not have.  The refinement site captures Python's
  control flow as a first-class mathematical object.

\item \textbf{Compositionality:} Product covers, incremental
  reverification, and the K\"unneth formula enable verification
  that scales beyond single functions to entire libraries.
\end{enumerate}

The cubical and cohomological structures are \emph{not redundant}.
The cubical interval parameterizes \emph{how} we prove (which proof
steps compose); the cohomological cover parameterizes \emph{what}
we prove (which inputs to handle).  The unification
(Theorem~\ref{thm:descent-hcomp}) shows they are aspects of a
single underlying structure: filling cubes from compatible
boundary data.

The compiler translates this theoretical framework into a practical
tool: proof terms become Z3 verification conditions, trust is
tracked as provenance, and annotations bind to code.  The result
is a verification system that is expressive enough for real Python
libraries, tractable enough for automation, and transparent enough
for human review.
